\subsection*{MT-M21. Identification with Meyer's zeta-summation quotient over \(\mathbb Q\)}
For the Riemann zeta function, the canonical quotient in MT-M18
is naturally isomorphic, as a topological convolution module,
to the zero-space in Meyer's classical spectral realization.
This identification is unconditional and includes all zero
multiplicities.

Here are the precise source and conventions. In Ralf Meyer,
\emph{A spectral interpretation for the zeros of the Riemann
zeta function}, G\"ottingen Seminars 2004--05, pp.~117--137,
let \(H_- =\mathcal O(\mathbb R_+^\times)\) be the strong
Schwartz space of Section~2, let \(H_+\) be the even Schwartz
space on \(\mathbb R\), and let
\[
 H_\cap=\{u\in H_+:u(0)=\widehat u(0)=0\},\qquad
 (Zu)(x)=\sum_{k\ge1}u(kx).
\]
His Theorem~3.3, printed p.~124 (PDF p.~8), proves the
closed-image assertion for the zeta summation map; Section~4,
printed pp.~125--126, defines
\(H_-^0=H_-/ZH_\cap\).
Theorem~4.1, printed p.~126 (PDF p.~10), characterizes its
Mellin-transform kernel by two half-plane division conditions.
Under the logarithmic coordinate map
\[
       \mathcal L:\mathcal E\longrightarrow H_-,
       (\mathcal Lf)(x)=f(\log x),
\]
one has the exact equality of closed subspaces
\[
          \mathcal L\bigl(\overline{\psi_\zeta*\mathcal E}\bigr)
                         = ZH_\cap.                  \tag{M21.1}
\]
Consequently \(Q_\zeta\simeq H_-^0\), and the adapted
absolute basis, nuclear convolution, and trace constructions
of MT-M18--MT-M20 apply to this independently defined
zeta-summation quotient with their stated positivity hypotheses.

\begin{proof}
The map \(\mathcal L\) is a topological isomorphism:
the invariant derivatives \((x\,d/dx)^j\) become
\(d^j/dt^j\) for \(t=\log x\), and requiring rapid
decay after every real power weight \(x^\sigma\)
is equivalent to the defining seminorms
\(\sup_t e^{N|t|}|f^{(j)}(t)|\). Polynomial factors in
\(t\), if included in the Schwartz definition, are
absorbed by increasing \(N\). Its Mellin transform is
exactly \(\mathcal Mf(s)\), since \(d^\times x=dt\).
Multiplicative convolution becomes additive convolution.

Let \(F=\mathcal Mf\). Meyer's Theorem~4.1 says that
\(\mathcal Lf\in ZH_\cap\) precisely when
\[
 t\longmapsto F(\sigma+it)/\zeta(\sigma+it)
          \text{ is Schwartz for every }\sigma\ge\tfrac12,
                                                               \tag{M21.2}
\]
and
\[
 t\longmapsto F(\sigma+it)/\zeta(1-\sigma-it)
          \text{ is Schwartz for every }\sigma\le\tfrac12.
                                                               \tag{M21.3}
\]
The conditions include holomorphic removal of the corresponding
zeros on those lines. If they hold, every nontrivial zero
\(\rho\) of real part at least \(1/2\) is a zero of \(F\)
to at least its full multiplicity, by (M21.2). For zeros in
the other half-plane, the zeta functional equation says that
\(\zeta(1-s)\) has the same nontrivial zero divisor as
\(\zeta(s)\); condition (M21.3) gives the same conclusion.
Thus \(f\) belongs to the full zero-jet kernel in MT-M18.

Conversely, suppose \(F\) vanishes on this entire divisor.
The function \(F/\zeta\) is holomorphic on
\(\operatorname{Re}s>-2\). All nontrivial zero singularities
are removed, the pole of \(\zeta\) at 1 gives a removable
zero of its reciprocal, and its first trivial zero is at
\(-2\), outside this open half-plane.

Fix \(A\ge2\). Use rectangles
\([-1,A+1]+i[h_n,h_{n+1}]\), with the horizontal levels
constructed in MT-M18 for \(\zeta\). The finite-\(L\)
Blaschke estimates in that proof give a fixed inverse
polynomial lower bound for \(|\zeta|\) on the horizontal
edges at large height. The two vertical edges are a fixed
positive distance from every nontrivial zero, so the same
disk argument gives their polynomial lower bound. They
also avoid all trivial zeros. The finitely many low-height
boundaries cause no problem: the levels \(h_n\) are never
zero, and both vertical lines are zero-free.

For every \(N\), rapid fixed-strip decay of \(F\)
therefore bounds \(|F/\zeta|\) on these boundaries by
\(C_{A,N,f}(2+|n|)^{-N}\), after increasing the decay
power used for \(F\). The maximum principle for the
holomorphic quotient gives the same bound inside each
rectangle. Since adjacent heights have uniformly bounded
gaps, \(F/\zeta\) decreases rapidly on the whole strip
\([-1,A+1]\). Cauchy's formula on small fixed disks
inside that strip gives rapid decay of all its vertical
derivatives on \([1/2,A]\). This proves (M21.2).

Apply this same argument to \(F^\vee(s)=F(1-s)\).
It is an entire strip-Schwartz function and vanishes on
the full nontrivial divisor, because that divisor is
invariant under \(s\mapsto1-s\). The result for
\(F^\vee/\zeta\), followed by \(s\mapsto1-s\),
is exactly (M21.3). Meyer's image criterion now gives
\(\mathcal Lf\in ZH_\cap\).

We have proved that \(\mathcal L^{-1}(ZH_\cap)\)
is precisely the full zero-jet kernel. MT-M18 identifies
that kernel with \(\overline{\psi_\zeta*\mathcal E}\),
proving (M21.1). Both are closed subspaces of the
topologically identified test-function spaces, so their
quotients, with quotient topology, are naturally
topologically isomorphic. The logarithmic identification
commutes with scaling and integrated convolution, giving
the stated module assertion.
\end{proof}

This settles a specified arithmetic comparison map, beyond
a zero-built abstract model. It is a comparison with
Meyer's Poisson/zeta-summation construction, not an
identification with an unspecified adelic quotient in
Masullo's draft or with an \(\ell\)-adic cohomology group.
The classical spectral realization itself is Meyer's
result; the proof here supplies the explicit closed-ideal
and stable-block crosswalk used in this investigation.
The full number-field idele-class construction is located
in Meyer, \emph{On a representation of the idele class
group related to primes and zeros of L-functions},
Duke Math. J.~127 (2005), 519--595, author version
\texttt{math/0311468v3}, Theorem~5.8 (pp.~43--44) and
Theorem~5.11 (p.~47). A sectorwise topological comparison
with that construction requires its precise image
criterion and is not inferred merely from equal spectra.
